% =============================================================================
% Акт патрулирования лесов
% Приказ Минприроды N 955 от 15.12.2021, Приложение 3
% Оформление: ГОСТ Р 7.0.97-2016
% Engine: LuaLaTeX
% =============================================================================
\documentclass[12pt, a4paper]{article}

% --- Язык ---
\usepackage[russian]{babel}

% --- Шрифты (LuaLaTeX + fontspec) ---
\usepackage{fontspec}
\defaultfontfeatures{Ligatures=TeX}
((% if font_path %))
\setmainfont{PTSerif}[
    Path           = {((* font_path *))},
    Extension      = {.ttf},
    UprightFont    = {*-Regular},
    BoldFont       = {*-Bold},
    ItalicFont     = {*-Italic},
    BoldItalicFont = {*-BoldItalic}
]
((% else %))
\setmainfont{PT Serif}
((% endif %))

% --- Поля по ГОСТ Р 7.0.97-2016 ---
\usepackage[
    a4paper,
    left=30mm,
    right=10mm,
    top=20mm,
    bottom=20mm
]{geometry}

% --- Типографика ---
\usepackage{setspace}
\setstretch{1.15}
\setlength{\parindent}{1.25cm}

% --- Пакеты ---
\usepackage{graphicx}
\usepackage{tabularx}
\usepackage{array}
\usepackage{booktabs}
\usepackage{fancyhdr}
\usepackage{lastpage}
\usepackage{microtype}
\usepackage{enumitem}
\usepackage{xcolor}
\usepackage{hyperref}
\hypersetup{unicode=true, colorlinks=false, pdfencoding=auto}

% --- Колонтитулы ---
\pagestyle{fancy}
\fancyhf{}
\fancyfoot[C]{\small Стр.~\thepage~из~\pageref{LastPage}}
\renewcommand{\headrulewidth}{0pt}
\renewcommand{\footrulewidth}{0.4pt}

% --- Без нумерации разделов ---
\setcounter{secnumdepth}{0}

% --- Линия для подписи ---
\newcommand{\sigline}[1]{%
    \noindent\begin{tabular}{p{0.45\textwidth} p{0.08\textwidth} p{0.40\textwidth}}
        #1 & & \rule{6cm}{0.4pt} \\
           & & {\small\hfill(подпись, расшифровка)\hfill} \\
    \end{tabular}%
}

% --- Поле «заполняется от руки» ---
\newcommand{\blankfield}[2][6cm]{%
    \noindent #2~\rule{#1}{0.4pt}%
}

% =============================================================================
\begin{document}
\thispagestyle{fancy}

% ===== ШАПКА =====
\begin{center}
    {\small\textsc{Министерство лесного хозяйства и охраны объектов животного мира}}\\
    {\small\textsc{Нижегородской области}}\\[4pt]
    {\normalsize\textbf{Варнавинское лесничество}}\\
((% if sub_district %))
    {\small ((* sub_district|e *)) участковое лесничество}\\
((% endif %))
\end{center}

\noindent\rule{\textwidth}{0.8pt}

\vspace{0.6cm}

% ===== НАЗВАНИЕ =====
\begin{center}
    {\Large\textbf{АКТ}}\\[2pt]
    {\large\textbf{патрулирования лесов}}\\[6pt]
    № \textbf{((* act_number|e *))} \hfill от «((* act_day|e *))»~((* act_month|e *))~((* act_year|e *))~г.
\end{center}

\vspace{0.4cm}

% ===== 1. СВЕДЕНИЯ О ПАТРУЛИРОВАНИИ =====
\section{1. Сведения о патрулировании}

\noindent
\begin{tabularx}{\textwidth}{@{}>{\bfseries}p{5.5cm}X@{}}
    Дата патрулирования: & ((* patrol_date|e *)) \\[3pt]
    Время начала: & ((* patrol_time_start|e *)) \\[3pt]
    Время окончания: & ((* patrol_time_end|e *)) \\[3pt]
    Способ патрулирования: & Автоматический акустический мониторинг (система Faun) \\[3pt]
    Маршрут / зона покрытия: & Зона акустических сенсоров, координаты центра: ((* lat *))°N, ((* lon *))°E \\
\end{tabularx}

% ===== 2. ЛИЧНЫЙ СОСТАВ =====
\section{2. Личный состав}

\noindent
\begin{tabularx}{\textwidth}{@{}>{\bfseries}p{5.5cm}X@{}}
((% if ranger_name %))
    Инспектор: & ((* ranger_name|e *)) \\[3pt]
((% else %))
    Инспектор: & \rule{8cm}{0.4pt} \\[3pt]
((% endif %))
((% if badge_number %))
    Удостоверение №: & ((* badge_number|e *)) \\[3pt]
((% else %))
    Удостоверение №: & \rule{5cm}{0.4pt} \\[3pt]
((% endif %))
    Орган лесной охраны: & Варнавинское лесничество \\
\end{tabularx}

% ===== 3. МЕСТО ОБСЛЕДОВАНИЯ =====
\section{3. Место обследования}

\noindent
\begin{tabularx}{\textwidth}{@{}>{\bfseries}p{5.5cm}X@{}}
    Субъект РФ: & Нижегородская область \\[3pt]
    Лесничество: & Варнавинское \\[3pt]
((% if sub_district %))
    Участковое лесничество: & ((* sub_district|e *)) \\[3pt]
((% else %))
    Участковое лесничество: & --- \\[3pt]
((% endif %))
((% if quarter %))
    Квартал: & ((* quarter|e *)) \\[3pt]
((% else %))
    Квартал: & --- \\[3pt]
((% endif %))
((% if compartment %))
    Выдел: & ((* compartment|e *)) \\[3pt]
((% else %))
    Выдел: & --- \\[3pt]
((% endif %))
    Координаты GPS: & ((* lat *))°N, ((* lon *))°E \\[3pt]
((% if forest_purpose %))
    Целевое назначение лесов: & ((* forest_purpose|e *)) \\
((% else %))
    Целевое назначение лесов: & требует уточнения \\
((% endif %))
\end{tabularx}

% ===== 4. РЕЗУЛЬТАТЫ ПАТРУЛИРОВАНИЯ =====
\section{4. Результаты патрулирования}

\subsection{4.1. Данные системы акустического мониторинга}

\noindent
\begin{tabularx}{\textwidth}{@{}>{\bfseries}p{5.5cm}X@{}}
    Тип акустического события: & ((* violation_type|e *)) \\[3pt]
    Уверенность классификатора: & ((* confidence *))\% \\[3pt]
    Уровень приоритета: & ((* gating_level|e *)) \\[3pt]
    Время обнаружения: & ((* detected_at|e *)) \\[3pt]
    ID алерта: & ((* incident_id|e *)) \\[3pt]
    Модель классификатора: & YAMNet v7 (FEATURE\_DIM=2048) \\
\end{tabularx}

\subsection{4.2. Предварительная правовая квалификация}

\noindent
Предполагаемая статья: \textbf{((* article|e *))}

((% if legal_articles %))
\vspace{0.3cm}
\noindent\textbf{Применимые правовые нормы (RAG-анализ):}

\noindent ((* legal_articles|e *))
((% endif %))

((% if drone_photo_path %))
\subsection{4.3. Данные БПЛА}

\noindent
\begin{center}
    \includegraphics[width=0.65\textwidth]{((* drone_photo_path *))}\\[4pt]
    {\small Снимок с БПЛА. Координаты: ((* lat *))°N, ((* lon *))°E}
\end{center}

((% if drone_comment %))
\vspace{0.2cm}
\noindent\textbf{Анализ снимка (Gemma 3 27B):}

\noindent ((* drone_comment|e *))
((% endif %))
((% endif %))

((% if ranger_photo_path or ranger_report %))
\subsection{4.4. Данные инспектора}

((% if ranger_photo_path %))
\noindent
\begin{center}
    \includegraphics[width=0.55\textwidth]{((* ranger_photo_path *))}\\[4pt]
    {\small Фото с места, сделанное инспектором}
\end{center}
((% endif %))

((% if ranger_report %))
\vspace{0.2cm}
\noindent\textbf{Описание обстановки:}

\noindent ((* ranger_report|e *))
((% endif %))
((% endif %))

% ===== 5. ПРИНЯТЫЕ МЕРЫ =====
\section{5. Принятые меры}

\noindent Автоматический алерт направлен инспектору через Telegram-бот <<Faun>>
в ((* detected_at|e *)).
((% if ranger_name %))
Инспектор ((* ranger_name|e *)) принял вызов и прибыл на место.
((% endif %))

% ===== 6. ТЕХНИЧЕСКИЕ СРЕДСТВА =====
\section{6. Применённые технические средства}

\begin{itemize}[nosep, leftmargin=1.5cm]
    \item Акустические сенсоры (система Faun) --- обнаружение и классификация звуковых событий
    \item Нейросеть YAMNet v7 (FEATURE\_DIM=2048) --- классификация акустических сигналов
    \item TDOA-триангуляция --- определение координат источника звука
    \item БПЛА (дрон) --- фотофиксация обстановки на месте
    \item Gemma 3 27B (Yandex Cloud AI Studio) --- анализ снимков с БПЛА
    \item YandexGPT --- юридизация описаний, правовая квалификация
\end{itemize}

% ===== 7. ПРИЛОЖЕНИЯ =====
\section{7. Приложения}

\begin{enumerate}[nosep, leftmargin=1.5cm]
((% if drone_photo_path %))
    \item Снимок с БПЛА (стр.~\pageref{LastPage} данного акта)
((% endif %))
((% if ranger_photo_path %))
    \item Фото с места, сделанное инспектором
((% endif %))
    \item Аудиозапись акустического события (хранится в системе Faun, ID: ((* incident_id|e *)))
    \item Спектрограмма акустического сигнала
\end{enumerate}

% ===== 8. ПОДПИСИ =====
\section{8. Подписи}

\vspace{0.6cm}

\sigline{Инспектор:}

\vspace{0.6cm}

\blankfield{Дата подписания:}

% ===== ФУТЕР =====
\vspace{1.5cm}
\noindent\rule{\textwidth}{0.4pt}

\begin{center}
    {\small\color{gray} Акт сформирован автоматически системой акустического мониторинга леса <<Faun>>}\\
    {\small\color{gray} (ForestGuard). Данные подлежат верификации инспектором лесной охраны.}\\
    {\small\color{gray} Основание: Приказ Минприроды России от 15.12.2021 № 955, п.~14, Приложение~3.}
\end{center}

\end{document}
